\documentclass[10pt,a4paper]{amsart}
\usepackage{amsmath,amssymb,amsthm,graphicx}
% Retained preamble for the Stacks Project chapter snapshot in this directory.
% The chapter's own preamble.tex is a copy of the shared Stacks Project preamble; the block
% below is sliced MECHANICALLY, line for line, from it. Nothing was added, removed or
% redefined; the wrapper class, the xr-hyper/hyperref wiring and the externaldocument list
% are omitted because the assembler supplies the document class.
%
% lines taken from preamble.tex: 3, 5, 9, 11, 12, 156, 157, 158, ... (51 lines total)
%
% package coverage restored from the chapter's own preamble (2026-10-13): xy, xy-2cell

\usepackage[all]{xy}
\xyoption{2cell}
\UseAllTwocells

\newenvironment{reference}{\comment}{\endcomment}
\newenvironment{slogan}{\comment}{\endcomment}
\newenvironment{history}{\comment}{\endcomment}
\theoremstyle{plain}
\newtheorem{theorem}[subsection]{Theorem}
\newtheorem{proposition}[subsection]{Proposition}
\newtheorem{lemma}[subsection]{Lemma}
\theoremstyle{definition}
\newtheorem{definition}[subsection]{Definition}
\newtheorem{example}[subsection]{Example}
\newtheorem{exercise}[subsection]{Exercise}
\newtheorem{situation}[subsection]{Situation}
\theoremstyle{remark}
\newtheorem{remark}[subsection]{Remark}
\newtheorem{remarks}[subsection]{Remarks}
\numberwithin{equation}{subsection}
\def\lim{\mathop{\mathrm{lim}}\nolimits}
\def\colim{\mathop{\mathrm{colim}}\nolimits}
\def\Spec{\mathop{\mathrm{Spec}}}
\def\Hom{\mathop{\mathrm{Hom}}\nolimits}
\def\Ext{\mathop{\mathrm{Ext}}\nolimits}
\def\SheafHom{\mathop{\mathcal{H}\!\mathit{om}}\nolimits}
\def\SheafExt{\mathop{\mathcal{E}\!\mathit{xt}}\nolimits}
\def\Sch{\mathit{Sch}}
\def\Mor{\mathop{\mathrm{Mor}}\nolimits}
\def\Ob{\mathop{\mathrm{Ob}}\nolimits}
\def\Sh{\mathop{\mathit{Sh}}\nolimits}
\def\NL{\mathop{N\!L}\nolimits}
\def\CH{\mathop{\mathrm{CH}}\nolimits}
\def\proetale{{pro\text{-}\acute{e}tale}}
\def\etale{{\acute{e}tale}}
\def\QCoh{\mathit{QCoh}}
\def\Ker{\mathop{\mathrm{Ker}}}
\def\Im{\mathop{\mathrm{Im}}}
\def\Coker{\mathop{\mathrm{Coker}}}
\def\Coim{\mathop{\mathrm{Coim}}}
\DeclareMathSymbol{\boxtimes}{\mathbin}{AMSa}{"02}
\def\QCohstack{\mathcal{QC}\!\mathit{oh}}
\def\Cohstack{\mathcal{C}\!\mathit{oh}}
\def\Spacesstack{\mathcal{S}\!\mathit{paces}}
\def\Quotfunctor{\mathrm{Quot}}
\def\Hilbfunctor{\mathrm{Hilb}}
\def\Curvesstack{\mathcal{C}\!\mathit{urves}}
\def\Polarizedstack{\mathcal{P}\!\mathit{olarized}}
\def\Complexesstack{\mathcal{C}\!\mathit{omplexes}}
\def\Pic{\mathop{\mathrm{Pic}}\nolimits}
\def\Picardstack{\mathcal{P}\!\mathit{ic}}
\def\Picardfunctor{\mathrm{Pic}}
\def\Deformationcategory{\mathcal{D}\!\mathit{ef}}
\usepackage{cleveref}
\crefname{originallemma}{original lemma}{original lemmas}
\Crefname{originallemma}{Original lemma}{Original lemmas}
\makeatletter
\newlabel{algebra-lemma-flat-fibre-smooth}{{00TF}{source}{}{sourcepointer-1}{}}
\newlabel{algebra-lemma-flat-fibre-smooth@cref}{{[originallemma][0][]00TF}{[1][0][]source}{}{}{}}
\makeatother
\crefname{originallemma}{original lemma}{original lemmas}
\Crefname{originallemma}{Original lemma}{Original lemmas}
\makeatletter
\newlabel{algebra-lemma-what-does-it-mean-again}{{0AS8}{source}{}{sourcepointer-2}{}}
\newlabel{algebra-lemma-what-does-it-mean-again@cref}{{[originallemma][0][]0AS8}{[1][0][]source}{}{}{}}
\makeatother
\crefname{originallemma}{original lemma}{original lemmas}
\Crefname{originallemma}{Original lemma}{Original lemmas}
\makeatletter
\newlabel{bootstrap-definition-property-transformation}{{03XZ}{source}{}{sourcepointer-3}{}}
\newlabel{bootstrap-definition-property-transformation@cref}{{[originallemma][0][]03XZ}{[1][0][]source}{}{}{}}
\makeatother
\crefname{originallemma}{original lemma}{original lemmas}
\Crefname{originallemma}{Original lemma}{Original lemmas}
\makeatletter
\newlabel{formal-spaces-lemma-diagonal-morphism-formal-algebraic-spaces}{{0AN2}{source}{}{sourcepointer-4}{}}
\newlabel{formal-spaces-lemma-diagonal-morphism-formal-algebraic-spaces@cref}{{[originallemma][0][]0AN2}{[1][0][]source}{}{}{}}
\makeatother
\crefname{originallemma}{original lemma}{original lemmas}
\Crefname{originallemma}{Original lemma}{Original lemmas}
\makeatletter
\newlabel{formal-spaces-lemma-structure-locally-noetherian}{{0GHM}{source}{}{sourcepointer-5}{}}
\newlabel{formal-spaces-lemma-structure-locally-noetherian@cref}{{[originallemma][0][]0GHM}{[1][0][]source}{}{}{}}
\makeatother
\crefname{originallemma}{original lemma}{original lemmas}
\Crefname{originallemma}{Original lemma}{Original lemmas}
\makeatletter
\newlabel{formal-spaces-section-adic}{{0AQ2}{source}{}{sourcepointer-6}{}}
\newlabel{formal-spaces-section-adic@cref}{{[originallemma][0][]0AQ2}{[1][0][]source}{}{}{}}
\makeatother
\crefname{originallemma}{original lemma}{original lemmas}
\Crefname{originallemma}{Original lemma}{Original lemmas}
\makeatletter
\newlabel{lemma-base-change-versal}{{0CXM}{source}{}{sourcepointer-7}{}}
\newlabel{lemma-base-change-versal@cref}{{[originallemma][0][]0CXM}{[1][0][]source}{}{}{}}
\makeatother
\crefname{originallemma}{original lemma}{original lemmas}
\Crefname{originallemma}{Original lemma}{Original lemmas}
\makeatletter
\newlabel{lemma-rs}{{0GI6}{source}{}{sourcepointer-8}{}}
\newlabel{lemma-rs@cref}{{[originallemma][0][]0GI6}{[1][0][]source}{}{}{}}
\makeatother
\crefname{originallemma}{original lemma}{original lemmas}
\Crefname{originallemma}{Original lemma}{Original lemmas}
\makeatletter
\newlabel{restricted-lemma-algebraize-rig-etale-affine}{{0AR4}{source}{}{sourcepointer-9}{}}
\newlabel{restricted-lemma-algebraize-rig-etale-affine@cref}{{[originallemma][0][]0AR4}{[1][0][]source}{}{}{}}
\makeatother
\crefname{originallemma}{original lemma}{original lemmas}
\Crefname{originallemma}{Original lemma}{Original lemmas}
\makeatletter
\newlabel{section-axioms-functors}{{07XZ}{source}{}{sourcepointer-10}{}}
\newlabel{section-axioms-functors@cref}{{[originallemma][0][]07XZ}{[1][0][]source}{}{}{}}
\makeatother
\title{Openness of the smooth locus of an adic versal morphism}
\author{The Stacks Project Authors}
\date{stacks-project:artin.tex:local-snapshot-20261012; GNU FDL 1.2 or later}
\begin{document}
\maketitle
\noindent\textit{Editorial note.} Statement span is lines 6329-6340 of sources/stacks-artin/pinned-original-source.tex; proof span is lines 6342-6449, the complete author proof (106 body lines).\par\medskip
\noindent\textit{Editorial note.} Navigation stubs: 10 cited-but-not-included labels are supplied as non-load-bearing pointers with their pinned tags where the snapshot tag index has them.\par\medskip
\noindent\textit{Editorial note.} extra\_packages is empty on purpose: the retained preamble already loads xy with options and a second declaration clashes in the assembler wrapper.\par\medskip
\noindent\textit{Editorial note.} Text is retained verbatim, including the author's own typos and placeholders.\par\medskip
\noindent\textit{Editorial note.} Queue index 2 of queue\_artin.json (pb=4897, stmt 6329-6340, proof 6342-6449).\par\medskip
\noindent\textit{Editorial note.} stacks-artin wave ar01 (2026-10-14): fresh adjudication accept per lane-s3/artin\_judgments.md; pool per Chief 10:5x direct dispatch (idle-first over S12's 01:5x note, SMC precedent); artin.tex 99 rows/67 with proof judged 46 accept / 21 reject / 5 AGP-overlap avoidance rows; patch-sd-envs contract spans scanned. Named omission: 'Some details omitted' for the etale descent of smoothness in the reduction to affine W - non-load-bearing; the Rees algebra key step is in-span.\par\medskip
\noindent\textit{Editorial note.} Cross-chapter same-label warning (NOT an occupancy conflict, Chief ruling 2026-10-13 evening): label 'lemma-openness-smoothness' also exists in stacks-snapshot-20261011 — a different proposition in a different chapter file.\par\medskip
\noindent\textit{Allowed standard prerequisites (editorial).}\begin{itemize}
\item Explicit prerequisite boundary. The cited inputs are prerequisites; the argument in the retained span is not replaced by any cited result.
\end{itemize}\medskip
\ifdefined\hypertarget\hypertarget{sourcepointer-1}{}\fi
\noindent\textit{Source reference.} Original lemma 00TF: sources/stacks-artin/pinned-original-source.tex, source lines 6449--6449. Cited by the retained statement or proof but not part of the collected spans; supplied only as a navigation pointer. The argument does not depend on its content.\par\medskip
\ifdefined\hypertarget\hypertarget{sourcepointer-2}{}\fi
\noindent\textit{Source reference.} Original lemma 0AS8: sources/stacks-artin/pinned-original-source.tex, source lines 6449--6449. Cited by the retained statement or proof but not part of the collected spans; supplied only as a navigation pointer. The argument does not depend on its content.\par\medskip
\ifdefined\hypertarget\hypertarget{sourcepointer-3}{}\fi
\noindent\textit{Source reference.} Original lemma 03XZ: sources/stacks-artin/pinned-original-source.tex, source lines 6449--6449. Cited by the retained statement or proof but not part of the collected spans; supplied only as a navigation pointer. The argument does not depend on its content.\par\medskip
\ifdefined\hypertarget\hypertarget{sourcepointer-4}{}\fi
\noindent\textit{Source reference.} Original lemma 0AN2: sources/stacks-artin/pinned-original-source.tex, source lines 6449--6449. Cited by the retained statement or proof but not part of the collected spans; supplied only as a navigation pointer. The argument does not depend on its content.\par\medskip
\ifdefined\hypertarget\hypertarget{sourcepointer-5}{}\fi
\noindent\textit{Source reference.} Original lemma 0GHM: sources/stacks-artin/pinned-original-source.tex, source lines 6449--6449. Cited by the retained statement or proof but not part of the collected spans; supplied only as a navigation pointer. The argument does not depend on its content.\par\medskip
\ifdefined\hypertarget\hypertarget{sourcepointer-6}{}\fi
\noindent\textit{Source reference.} Original lemma 0AQ2: sources/stacks-artin/pinned-original-source.tex, source lines 6449--6449. Cited by the retained statement or proof but not part of the collected spans; supplied only as a navigation pointer. The argument does not depend on its content.\par\medskip
\ifdefined\hypertarget\hypertarget{sourcepointer-7}{}\fi
\noindent\textit{Source reference.} Original lemma 0CXM: sources/stacks-artin/pinned-original-source.tex, source lines 2013--2022. Cited by the retained statement or proof but not part of the collected spans; supplied only as a navigation pointer. The argument does not depend on its content.\par\medskip
\ifdefined\hypertarget\hypertarget{sourcepointer-8}{}\fi
\noindent\textit{Source reference.} Original lemma 0GI6: sources/stacks-artin/pinned-original-source.tex, source lines 6175--6178. Cited by the retained statement or proof but not part of the collected spans; supplied only as a navigation pointer. The argument does not depend on its content.\par\medskip
\ifdefined\hypertarget\hypertarget{sourcepointer-9}{}\fi
\noindent\textit{Source reference.} Original lemma 0AR4: sources/stacks-artin/pinned-original-source.tex, source lines 6449--6449. Cited by the retained statement or proof but not part of the collected spans; supplied only as a navigation pointer. The argument does not depend on its content.\par\medskip
\ifdefined\hypertarget\hypertarget{sourcepointer-10}{}\fi
\noindent\textit{Source reference.} Original lemma 07XZ: sources/stacks-artin/pinned-original-source.tex, source lines 2501--2619. Cited by the retained statement or proof but not part of the collected spans; supplied only as a navigation pointer. The argument does not depend on its content.\par\medskip
\begin{lemma}
\label{lemma-openness-smoothness}
Let $S$ be a locally Noetherian scheme. Let $V$ be a scheme locally
of finite type over $S$. Let $Z \subset V$ be closed. Let $W$ be
a locally Noetherian formal algebraic space over $S$ such that
$W_{red}$ is locally of finite type over $S$. Let $g : V_{/Z} \to W$
be an adic morphism of formal algebraic spaces over $S$. Let $v \in V$
be a closed point such that $g$ is versal at $v$ (as in
Section \ref{section-axioms-functors}).
Then after replacing $V$ by an open neighbourhood of $v$ the
morphism $g$ is smooth (see proof).
\end{lemma}

\begin{proof}
Since $g$ is adic it is representable by algebraic spaces (Formal Spaces,
Section \ref{formal-spaces-section-adic}).
Thus by saying $g$ is smooth we mean that $g$ should be smooth
in the sense of
Bootstrap, Definition \ref{bootstrap-definition-property-transformation}.

\medskip\noindent
Write $W = \colim W_n$ as in Formal Spaces, Lemma
\ref{formal-spaces-lemma-structure-locally-noetherian}.
Set $V_n = V_{/Z} \times_{\hat x, W} W_n$.
Then $V_n$ is a closed subscheme with underlying set $Z$.
Smoothness of $V \to W$ is equivalent to the smoothness
of all the morphisms $V_n \to W_n$ (this holds because any morphism
$T \to W$ with $T$ a quasi-compact scheme factors through
$W_n$ for some $n$). We know that the morphism $V_n \to W_n$
is smooth at $v$ by
Lemma \ref{lemma-base-change-versal}\footnote{The lemma applies since
the diagonal of $W$ is representable by algebraic spaces and
locally of finite type, see Formal Spaces, Lemma
\ref{formal-spaces-lemma-diagonal-morphism-formal-algebraic-spaces}
and we have seen that $W$ has (RS) in the proof of Lemma \ref{lemma-rs}.}.
Of course this means that given any $n$ we can shrink $V$
such that $V_n \to W_n$ is smooth. The problem is to find
an open which works for all $n$ at the same time.

\medskip\noindent
The question is local on $V$, hence we may assume $S = \Spec(R)$ and
$V = \Spec(A)$ are affine.

\medskip\noindent
In this paragraph we reduce to the case where $W$ is an affine formal
algebraic space. Choose an affine formal scheme $W'$ and an \'etale morphism
$W' \to W$ such that the image of $v$ in $W_{red}$ is in the
image of $W'_{red} \to W_{red}$. Then $V_{/Z} \times_{g, W} W' \to V_{/Z}$
is an adic \'etale morphism of formal algebraic spaces over $S$
and $V_{/Z} \times_{g, W} W'$ is an affine formal algebraic space.
By Algebraization of Formal Spaces,
Lemma \ref{restricted-lemma-algebraize-rig-etale-affine}
there exists an \'etale morphism $\varphi : V' \to V$ of affine schemes
such that the completion of $V'$ along $Z' = \varphi^{-1}(Z)$
is isomorphic to $V_{/Z} \times_{g, W} W'$ over $V_{/Z}$.
Observe that $v$ is the image of some $v' \in V'$.
Since smoothness is preserved under base change we see that
$V'_n \to W'_n$ is smooth for all $n$. In the next paragraph
we show that after replacing $V'$ by an open neighbourhood of $v'$
the morphisms $V'_n \to W'_n$ are smooth for all $n$.
Then, after we replace $V$ by the open image of $V' \to V$,
we obtain that $V_n \to W_n$ is smooth by \'etale descent of smoothness.
Some details omitted.

\medskip\noindent
Assume $S = \Spec(R)$, $V = \Spec(A)$, $Z = V(I)$, and $W = \text{Spf}(B)$.
Let $v$ correspond to the maximal ideal $I \subset \mathfrak m \subset A$.
We are given an adic continuous $R$-algebra homomorphism
$$
B \longrightarrow A^\wedge
$$
Let $\mathfrak b \subset B$ be the ideal of topologically nilpotent
elements (this is the maximal ideal of definition of the Noetherian adic
topological ring $B$). Observe that $\mathfrak b A^\wedge$ and
$IA^\wedge$ are both ideals of definition of the Noetherian adic
ring $A^\wedge$. Also, $\mathfrak m A^\wedge$ is a maximal ideal
of $A^\wedge$ containing both $\mathfrak b A^\wedge$ and $IA^\wedge$.
We are given that
$$
B_n = B/\mathfrak b^n \to A^\wedge/\mathfrak b^n A^\wedge = A_n
$$
is smooth at $\mathfrak m$ for all $n$. By the discussion above
we may and do assume that $B_1 \to A_1$ is a smooth ring map.
Denote $\mathfrak m_1 \subset A_1$ the maximal ideal corresponding
to $\mathfrak m$. Since smoothness implies flatness, we see that:
for all $n \geq 1$ the map
$$
\mathfrak b^n/\mathfrak b^{n + 1} \otimes_{B_1} (A_1)_{\mathfrak m_1}
\longrightarrow
\left(\mathfrak b^nA^\wedge/\mathfrak b^{n + 1}A^\wedge\right)_{\mathfrak m_1}
$$
is an isomorphism (see
Algebra, Lemma \ref{algebra-lemma-what-does-it-mean-again}).
Consider the Rees algebra
$$
B' = \bigoplus\nolimits_{n \geq 0} \mathfrak b^n/\mathfrak b^{n + 1}
$$
which is a finite type graded algebra over the Noetherian ring $B_1$ and
the Rees algebra
$$
A' = \bigoplus\nolimits_{n \geq 0}
\mathfrak b^nA^\wedge/\mathfrak b^{n + 1}A^\wedge
$$
which is a a finite type graded algebra over the Noetherian ring $A_1$.
Consider the homomorphism of graded $A_1$-algebras
$$
\Psi : B' \otimes_{B_1} A_1 \longrightarrow A'
$$
By the above this map is an isomorphism after localizing at
the maximal ideal $\mathfrak m_1$ of $A_1$.
Hence $\Ker(\Psi)$, resp.\ $\Coker(\Psi)$ is a finite module
over $B' \otimes_{B_1} A_1$, resp.\ $A'$ whose localization
at $\mathfrak m_1$ is zero. It follows that after replacing
$A_1$ (and correspondingly $A$) by a principal localization
we may assume $\Psi$ is an isomorphism. (This is the key step of the proof.)
Then working backwards we see that $B_n \to A_n$ is flat, see
Algebra, Lemma \ref{algebra-lemma-what-does-it-mean-again}.
Hence $A_n \to B_n$ is smooth (as a flat ring map with smooth
fibres, see Algebra, Lemma \ref{algebra-lemma-flat-fibre-smooth})
and the proof is complete.
\end{proof}

\end{document}
